\subsection{Unital algebras}

A unital algebra over K is a pair \((A,e)\) where A is an algebra over K with identity element, and e is the identity element of A. Since e is uniquely determined by A, we will sometimes say, by abuse of language, that A is a unital algebra. If \((A,e)\) and \((A',e')\) are unital algebras, a unital morphism of \((A,e)\) into \((A',e')\) is a morphism \(\varphi\) from A to \(A'\) such that \(\varphi(e)=e'\). A unital subalgebra of \((A,e)\) is a pair \((A',e)\) , where \(A'\) is a subalgebra of A containing e.

We will often denote the identity element by 1.

Lemma 1. --- Let A be an algebra. For every idempotent j of A, the subspace jAj of A is the set of \(x \in A\) such that xj = jx = x. This is a subalgebra of A which admits the unit element j.

The proof is elementary.

